\clearpage
\section{Stasjonær sensorstudie 23.10.2026}\label{sec:sensorstudie23}
Økten prioriterte sensorstabilitet fra revisjon \kode{02b542f}, etter
sikkerhetskopiering. De to nye måleseriene nedenfor supplerer de
tidligere kjøreobservasjonene i avsnitt~\ref{sec:utvikling23}.

\subsection{Oppsett, hypoteser og kontrollert innsamling}
Bruker bekreftet at motorstrømmen var frakoblet, mens USB og sensorer var
strømforsynt. En separat Ada-applikasjon initialiserer motorutgangen med
stopp og sentrerer servoen til 84\textdegree{}. Den inneholder ingen kall
til kjørefunksjonen. Etter siste måling sendes stopp igjen, og programmet
venter uten nye bevegelseskommandoer. Dette er en fysisk sensortest,
ikke en test under motorbelastning.

\begin{table}[htbp]
\centering\small
\caption{Pinnetilordning i diagnostikken og dagens navigasjonskode.}
\begin{tabular}{lll}\toprule
Sensor & TRIG & ECHO \\\midrule
Front & P13 & P1 \\
Venstre & P14 & P2 \\
Høyre & P15 & P8 \\\bottomrule
\end{tabular}
\end{table}
Bruker oppga to motstander på 1 k\(\Omega\) og 1,2 k\(\Omega\) per
ECHO-spenningsdeler. Plasseringen er ikke fysisk bekreftet. Et gjenfunnet
eldre diagram gjelder én sensor og andre motstandsverdier; det er derfor
ikke brukt som bekreftelse på dagens krets. Diagrammet brukte dessuten
P2 og P14 til LED. Bruker bekreftet etter forsøkene at bare venstre
ultralydsensor bruker disse pinnene. Hypotesen om delt LED-last legges
derfor til side; motstandsrekkefølgen er fortsatt uavklart.

Databladet anbefaler mer enn 60 ms mellom målesykluser\cite{hcsr04}.
Diagnostikken utløser bare én sensor om gangen. Hver serie har følgende
faste rekkefølge: 40 runder front--venstre--høyre med 65 ms mellom globale
utløsninger, 40 tilsvarende runder med 100 ms, deretter 40 målinger på
venstre alene med hvert intervall. Totalt blir dette 320 avlesninger per
serie. En bestemt sensor måles dermed sjeldnere enn det globale
triggerintervallet når alle tre er aktive.

Serie 1 brukte 3 ms frist for ECHO-stigning. Tidsdata viste at gyldige
venstreekko nærmet seg denne fristen. Serie 2 økte derfor bare denne
diagnostikkfristen til 10 ms. Begge hadde 23 ms frist for at en påbegynt
puls skulle avsluttes. Ved oppstart av hver avlesning kontrolleres også
om ECHO allerede er høyt. Feilkodene skiller manglende stigning,
manglende fall, allerede høyt signal og avstand utenfor 2--200 cm.
En timeout gir aldri en brukbar fri avstand.

\begin{figure}[htbp]
\centering
\begin{tikzpicture}[node distance=5mm,
 s/.style={draw=usnpurple,rounded corners,fill=softpurple,
 text width=10cm,align=center,font=\small,inner sep=5pt}]
\node[s] (wait) {Vent til planlagt globalt triggertidspunkt};
\node[s,below=of wait] (read) {Kontroller ECHO, utløs én sensor,\par
 vent på stigning og fall innen eksplisitte frister};
\node[s,below=of read] (log) {Lagre status, puls, lesetid og sekvensnummer};
\draw[-{Latex},usnpurple] (wait)--(read);
\draw[-{Latex},usnpurple] (read)--(log);
\end{tikzpicture}
\caption{Sekvensiell, stasjonær måleprotokoll. Motorene holdes stoppet.}
\end{figure}

\subsection{Modeller og måleusikkerhet}
Nominell avstand fra pulsbredde er
\begin{equation}
 d_{\mathrm{cm}}=\frac{t_{\mathrm{puls},\mu s}}{58}.
\end{equation}
Dette er databladets omregningsmodell\cite{hcsr04}, ikke dokumentert
nøyaktighet for det aktuelle oppsettet. Den nye rene Ada-modulen
\kode{Adaptive_Range_Model} implementerer omregningen, gyldig
arbeidsområde og spenningsdelerberegningen. Diagnostikken lagrer både
klokkebasert pulsbredde og den opprinnelige løkketelleren.

Inspeksjon av runtime-filen \kode{s-bbbosu.adb} viser at
\kode{Ada.Real_Time.Clock} bygger på RTC ved 32\,768 Hz. Ett klokketrinn
er derfor omtrent 30,5 \(\mu\)s, tilsvarende 0,526 cm i modellen.
Polling av ECHO, avrunding og avbrudd kommer i tillegg. Mikrosekunitsnavnet
i loggen betyr ikke oppløsning på ett mikrosekund. Den eldre
\kode{Delay_Us}-rutinen bruker en håndjustert NOP-løkke. Forskjellen
mellom de to avstandsestimatene kan undersøkes, men ingen av dem er
kalibrert mot kjente avstander i disse seriene.

For spenningsdeleren gjelder
\begin{equation}
 V_{\mathrm{signal}}=V_{\mathrm{ECHO}}
 \frac{R_{\mathrm{GND}}}{R_{\mathrm{serie}}+R_{\mathrm{GND}}}.
\end{equation}
Ved antatt ECHO på 5 V gir 1 k\(\Omega\) i serie og 1,2 k\(\Omega\)
til GND omtrent 2,73 V. Motsatt plassering gir omtrent 2,27 V.
Micro:bit dokumenterer \(V_{IH,\min}=0{,}7V_{DD}\)\cite{pins}; ved
antatt \(V_{DD}=3{,}3\) V blir dette 2,31 V. Den motsatte plasseringen
kan dermed mangle garantert margin for logisk høy. Dette regnestykket
bekrefter verken dagens motstandsrekkefølge, faktisk spenning eller
årsaken til venstresensorens feil. Resistortoleranse og forsyningsvariasjon
må også undersøkes.

For gyldige målinger beregner analyseverktøyet gjennomsnitt, median og
utvalgsstandardavvik. Med en faktisk oppmålt referanse \(\hat d_i\)
beregnes også
\begin{align}
 \operatorname{bias}&=\frac1n\sum_{i=1}^{n}(d_i-\hat d_i),\\
 \operatorname{MAE}&=\frac1n\sum_{i=1}^{n}|d_i-\hat d_i|,\\
 \operatorname{RMSE}&=\sqrt{\frac1n\sum_{i=1}^{n}(d_i-\hat d_i)^2}.
\end{align}
Referanseavstander ble ikke oppgitt for disse seriene. Feilmetrikkene
lagres derfor som \kode{null}. Gjennomsnittet brukes ikke som en oppdiktet
fasit. Ugyldige målinger inngår i feilandelens nevner, men aldri som
0 cm i gjennomsnittet.

\subsection{Resultater fra to fysiske serier}
Begge seriene fullførte med 320 sammenhengende sekvensnumre og
sluttmelding om stopp. Av 640 avlesninger var
front og høyre gyldige i alle sine 160 forsøk hver. Venstre ga 36 av 160
gyldige avlesninger i serie 1 og 115 av 160 i serie 2.

\begin{table}[htbp]
\centering\small
\caption{Venstre kanal, 40 avlesninger per fase. Gyldighet er
diagnostikkens status, ikke bekreftet avstandsnøyaktighet.}
\begin{tabular}{llrrrr}\toprule
Serie & Fase & Gyldige & Ugyldige & Andel ugyldige & \(s\), cm \\\midrule
1 & Alle, 65 ms &16&24&60,0\%&4,52\\
1 & Alle, 100 ms&10&30&75,0\%&5,66\\
1 & Bare venstre, 65 ms&5&35&87,5\%&2,90\\
1 & Bare venstre, 100 ms&5&35&87,5\%&2,13\\\midrule
2 & Alle, 65 ms&29&11&27,5\%&24,15\\
2 & Alle, 100 ms&28&12&30,0\%&23,06\\
2 & Bare venstre, 65 ms&30&10&25,0\%&25,88\\
2 & Bare venstre, 100 ms&28&12&30,0\%&23,54\\\bottomrule
\end{tabular}
\label{tab:left-sensor23}
\end{table}

\begin{table}[htbp]
\centering\small
\caption{Deskriptive avstander i fasen med alle sensorer og 65 ms
globalt intervall. Hver rad har 40 forsøk; spredning gjelder gyldige ekko.}
\begin{tabular}{llrrrr}\toprule
Serie & Sensor & Gyldige & Snitt, cm & Median, cm & \(s\), cm \\\midrule
1&Front&40&80,47&82,09&5,37\\
1&Venstre&16&82,08&82,35&4,52\\
1&Høyre&40&70,03&69,98&0,56\\
2&Front&40&80,31&82,60&6,35\\
2&Venstre&29&54,40&55,78&24,15\\
2&Høyre&40&70,02&69,98&0,14\\\bottomrule
\end{tabular}
\end{table}

I serie 1 var alle 124 ugyldige venstremålinger manglende ECHO-stigning.
Serie 2 hadde 44 slike feil og én avstand utenfor arbeidsområdet.
Maksimal observert lesetid var 7935 \(\mu\)s i serie 1 og
10\,071 \(\mu\)s i serie 2. Disse er observerte verdier fra en avgrenset
stasjonær test, ikke beviste WCET-grenser eller maksimal
hindring-til-motorstopp-responstid.

\subsection{Filterforsøk og vurdering av årsak}
To analysekandidater bruker ikke-overlappende tripler per fase og sensor.
Medianvarianten krever tre gyldige ekko. Den konsistente
minimumsvarianten krever i tillegg høyst 15 cm spenn og beholder den
nærmeste avstanden. Begge avviser en trippel med ugyldig måling.
Én restmåling per gruppe på 40 inngår ikke i trippelanalysen.
Serie 1 ga ingen aksepterte venstretripler for noen kandidat. Et filter
kan altså ikke gjenopprette pålitelig avstand fra disse hullene.
Kandidatene erstatter ikke den aktive navigasjonens filter.

Økt stigningsfrist ga flere mottatte ekko, men også langt større
avstandsspredning på venstre. Dette er ikke tilstrekkelig grunnlag for
å ta 10 ms inn som en ferdig retting i navigasjonen. Isolert måling
eller 100 ms intervall fjernet heller ikke feilen i disse seriene.
Fast faserekkefølge og manglende uavhengig kontroll av målscenen
begrenser årsaksslutninger. Dårlig spenningsmargin, kontaktfeil,
modulvariant og refleksjoner står fortsatt som
hypoteser. Motorstøy er ikke undersøkt med motorene i drift.

Neste forsøk bør begynne med en oppmålt, stor målflate 60 cm fra
venstre sensor og fysisk bekreftelse av
spenningsdeleren. Deretter kan venstre/høyre sensormodul byttes med
strømmen av og bare én variabel endres om gangen. Måleserier ved
20, 30, 45, 60, 100 og 150 cm er protokollført som videre arbeid;
de er ikke gjennomført eller presentert som resultater.

\subsection{Programvarekontroll, sporbarhet og sanntidsteori}
De 1729 eksisterende Ada-kontrollene besto. En ny vertsbasert Ada-suite
bestod 24 kontroller av pulsmodell, grenseverdier, ugyldige inndata og
betinget spenningsmargin. Analyseverktøyet besto 16 Python-enhetstester
med syntetiske data, blant annet kjente MAE/RMSE-verdier,
feilinjeksjon i CSV, sekvenskontroll og ugyldighetsveto. Antallet
Ada-kontroller er ikke antall fysiske forsøk eller dekningsprosent.
Både navigasjonsfirmware og separat diagnostikk ble bygget.

Sanntidsfaget definerer responstid som tiden fra inngang til nødvendig
utgang, og beskriver fail-safe som bevart integritet med midlertidig
stans\cite[PDF-s.~20 og 74]{forelesninger2026}. En sensorrutines
lesetid alene dekker derfor ikke robotens reaksjonstid. Protokollen
opprettholder sikker stopp når nødvendig informasjon er ukjent.
Ingen nye påstander om fysisk svingklarering, stopplengde, PID,
IMU-fusjon eller full autonom gulvkjøring følger av denne studien.

Rapportprosjektet inneholder begge rådatafiler (\kode{sensor-run-01.csv}
og \kode{sensor-run-02.csv}), tilhørende JSON-analyser og
\kode{SENSORPROTOKOLL.md}. Git-repoets evidensmappe bevarer dessuten
seriellogger, firmwarehash, kildeøyeblikksbilder og bygg-/testlogger.
